% title: 站点首页
% summary: 一个把 LaTeX 内容直接编译成 HTML 的数学博客骨架，适合课程笔记、研究随笔和学术长文发布。

\section{这个站点现在能做什么}

这个版本不再沿用旧站的 Markdown 渲染流程，而是把内容源统一切换为 \textbf{LaTeX}。也就是说，后续你只需要在 `PostCode` 目录下维护 `.tex` 文件，`WebCode/build.py` 就会把它们编译成结构化的 HTML 页面。

\tableofcontents

\subsection{目录分工}

\begin{enumerate}
\item `WebCode` 保存站点框架、模板、样式、脚本，以及 LaTeX 到 HTML 的 Python 构建程序。
\item `PostCode` 保存内容本身，包括首页、独立页面、文章以及图片等资源。
\item 构建输出会进入 `WebCode/dist`，方便你后续直接部署到 GitHub Pages 或其他静态托管平台。
\end{enumerate}

\subsection{这条新链路为什么更适合你}

\begin{definition}
这里的“更适合”是指：内容创作习惯、数学公式表达、长文结构控制三件事，都尽量围绕 LaTeX 本身展开，而不是先写 Markdown 再做兼容修补。
\end{definition}

对于数学站点而言，行内公式 $f(x)=x^2+\sin x$、块级公式
\[
\int_0^1 x^2\,dx = \frac{1}{3}
\]
以及像 `align` 这样的环境，本来就是 LaTeX 最自然的表达方式。把内容层直接建立在 LaTeX 上，后续维护会轻松很多。

\subsection{建议你的下一步}

\begin{itemize}
\item 把你原来准备发布的内容逐步迁移到 `PostCode/posts/` 里。
\item 如果文章较长，可以按专题建立子目录，例如 `PostCode/posts/analysis/`、`PostCode/posts/algebra/`。
\item 如果需要插图，统一放到 `PostCode/assets/` 中，再用 `\includegraphics` 调用。
\end{itemize}

\begin{remark}
这一版已经足够支撑“课程笔记型”和“研究笔记型”网站。如果你接下来想加自动引用、文献系统、搜索页、标签页，我们可以继续往上接。
\end{remark}
